% TOP_PROBLEM: {"id": 4700023, "title": "An extended Poncelet problem I", "category_id": 6, "category": {"id": 6, "name": "geometry", "display_name": "Geometry", "description": "Euclidean and non-Euclidean geometry, geometric structures.", "slug": "geometry", "order_index": 6, "created_at": "2026-07-31T15:26:25.671Z"}, "status": "open"}
% FIRST_POSTED: null
% ENTRY_KIND: statement_only
% Imported from: batch15.tex; UnsolvedMath ID 4700023
\documentclass[11pt]{article}
\usepackage[margin=25mm]{geometry}
\usepackage{fontspec}
\setmainfont{Latin Modern Roman}
\usepackage{amsmath,amssymb,mathtools}
\usepackage{unicode-math}
\setmathfont{Latin Modern Math}
\usepackage{xurl,hyperref}
\hypersetup{colorlinks=true,urlcolor=blue}
\setcounter{secnumdepth}{0}
\setlength{\parindent}{0pt}
\setlength{\parskip}{5pt}
\setlength{\emergencystretch}{3em}
\newcommand{\sref}[2]{\hyperlink{source:#1:#2}{[S#2]}}
\newcommand{\cataloguescope}[1]{\paragraph{Recorded target.}#1}
\begin{document}
% UnsolvedMath ID 4700023; source catalogue code AMR-046-0023
\setcounter{section}{4700022}
\section{An extended Poncelet problem I}\label{4700023}
{\small\sffamily UnsolvedMath ID 4700023 · Geometry}
\subsection{Definitions and mathematical statement}

\paragraph{Shared notation.}$\mathbb N=\{1,2,\ldots\}$, and $\mathbb Z,\mathbb Q,\mathbb R,\mathbb C$ denote the integers, rationals, reals and complex numbers. Any other conventions are specified locally.


An algebraic oval is a smooth compact connected component of a real plane algebraic curve, homeomorphic to a circle. Choose convex ovals $\gamma,\Gamma$ with $\gamma$ strictly inside the region bounded by $\Gamma$. The Poncelet map $P:\Gamma\to\Gamma$ follows a tangent line to the inner oval from $p\in\Gamma$ to its other intersection with $\Gamma$, consistently choosing the first of the two tangent intersections encountered in the counterclockwise direction. Only pairs for which this procedure defines a circle homeomorphism are considered. Rotation conjugacy means there are a homeomorphism $h:\Gamma\to\mathbb R/\mathbb Z$ and $\rho\in\mathbb R/\mathbb Z$ with $h(P(p))=h(p)+\rho$ for all $p$; no rationality condition on $\rho$ is imposed. A curve of degree $d$ is given by an irreducible real polynomial of total degree $d$. The chosen oval is part of that curve; the curve may have other real components.
\paragraph{Statement recorded in the supplied source.}
Are there two such irreducible curves of degrees $n,m$ with $n+m>4$, each carrying a nested convex oval as above, whose well-defined Poncelet map is conjugate to a circle rotation? \sref{4700023}{2}
\subsection{Short English statement}
Can nested ovals of irreducible curves beyond the two-conic case have a rotation-conjugate Poncelet map?
\subsection{Sources}
\begin{description}
\item[\hypertarget{source:4700023:1}{S1}] ULAM.AI and the UnsolvedMath Contributors, \emph{UnsolvedMath: A Curated Collection of Open Mathematics Problems}, record 4700023 (AMR-046-0023). CC BY 4.0. \url{https://huggingface.co/datasets/ulamai/UnsolvedMath}.
\item[\hypertarget{source:4700023:2}{S2}] A. Gasull, Some open problems in low dimensional dynamical systems (2020), Section 6.2, definition of the Poncelet map, Theorem 6.2 and full Problem 23, pp.23--24. \url{https://arxiv.org/pdf/2012.02524}.
\end{description}
\label{end:4700023}
\end{document}
